\section{Related work}\label{sec:related}

\subsection{Learner models and knowledge tracing}

Knowledge tracing estimates a learner's evolving knowledge from response sequences. Bayesian knowledge tracing models each knowledge component as a hidden binary state \citep{corbett1994bkt}; logistic models such as the additive factor model (AFM) and performance factors analysis (PFA) combine a learner ability, item difficulty and practice counts per knowledge component \citep{cen2006lfa,pavlik2009pfa,pelanek2017learner}, in the sense of knowledge components described by \citet{koedinger2012kli}. Recurrent and attention-based models, starting with DKT \citep{piech2015dkt} and continuing with SAKT and AKT \citep{pandey2019sakt,ghosh2020akt}, learn the dynamics from data \citep[see][for a survey]{abdelrahman2023survey}. Several comparisons found that well-specified logistic models match or approach deep models in forecasting accuracy \citep{khajah2016deep,wilson2016basics,gervet2020deep}. DKT's step-to-step predictions can also move in implausible directions, for instance lowering the predicted mastery of a skill after a correct answer \citep{yeung2018dktplus}, which has motivated consistency and monotonicity constraints \citep{zhang2023cmkt}. For programming, KT has been extended with code features \citep{shi2022codedkt}, and adaptive navigation support has long guided learners to programming problems of suitable difficulty \citep{hsiao2010guiding,brusilovsky2007user}. Evaluation of learner models is nonetheless dominated by held-out predictive accuracy, although frameworks for evaluating them by their consequences for tutoring decisions have been proposed \citep{gonzalezbrenes2015leopard,huang2015multifaceted}.

\subsection{Reinforcement learning for instructional sequencing and exercise recommendation}

Instructional sequencing with RL has a long history. Policies induced from data have been tested with students in tutoring systems \citep{chi2011umuai,shen2018jedm,zhou2022hrl}, in online courses \citep{bassen2020chi} and in a narrative mathematics tutor \citep{ruan2024rltutor}; related work has used POMDP planning \citep{rafferty2016pomdp}, multi-armed bandits \citep{clement2015bandits} and approximately optimal teaching models \citep{whitehill2018teaching}. Reviews conclude that successes are real but uneven: \citet{doroudi2019reward} found that policies derived from data-driven student models often failed to outperform reasonable baselines when tested with students and that success was more likely when policies were grounded in theory-driven models, and more recent overviews point to methodological issues and to the need for larger evaluations with actual learners \citep{singla2021rleducation,riedmann2025rlreview}. In parallel, the recommender-systems literature has built deep RL recommenders of exercises and learning paths that rely on a fitted learner model---usually a KT model---for the state, the reward or the measure of learning: multi-objective exercise recommendation \citep{huang2019dre}, cognitive-structure-aware learning paths \citep{liu2019cseal}, concept-aware DKT recommendation \citep{ai2019cdkt}, set-to-sequence and hierarchical path planners \citep{chen2023src,li2023gehrl}, knowledge-state distillation \citep{li2024pksd} and ExRec, which defines rewards and task scores through calibrated KT states \citep{ozyurt2025exrec}. The UniER benchmark unifies item- and path-level recommenders under a KT-derived weighted cognitive gain \citep{cheng2026unier}, and deep RL learning-path systems also appear in this journal \citep{ruan2025adaptive}. These works advance representation and optimisation; our question is how much of a KT-measured gain survives when the world that generates outcomes is not the model that measures them.

\subsection{From predictive accuracy to instructional decisions}

That forecasting accuracy does not settle instructional decisions is a recurring finding. For a substantial share of students, population knowledge-tracing parameters implied about twice as much practice as individually fitted parameters \citep{lee2012individualizing}; student models of similar predictive accuracy can imply substantially different amounts of practice \citep{rollinson2015predictive}; knowledge-tracing parameters can be semantically degenerate even when they fit well \citep{doroudi2017identifiability}; and misspecified models can make mastery decisions inequitable across learners \citep{doroudi2019fairer}. \citet{gonzalezbrenes2015leopard} argued for evaluating student models by the tutoring decisions they induce. The robust evaluation matrix \citep{doroudi2017rem} evaluates each candidate policy under several student models and cautions against trusting a policy that looks good only under the model used to derive it. Off-policy estimators based on importance sampling evaluate policies from logged data without a simulator \citep{mandel2014offline,thomas2016dataefficient}, but they require logged action probabilities, which self-directed practice on an online judge does not provide, so simulation remains the default for exercise recommenders. In machine learning more broadly, predictors with equivalent held-out accuracy can behave very differently downstream \citep{breiman2001two,damour2022underspecification}.

\subsection{Optimising against learned models and proxy rewards}

In model-based RL, a policy optimised in a learned dynamics model exploits the model's errors, especially where data are scarce; ensembles, short model rollouts and uncertainty penalties were introduced to limit this \citep{kurutach2018metrpo,janner2019trust,yu2020mopo,moerland2023mbrl,levine2020offline}. Optimising a proxy objective can likewise raise the proxy while the intended objective stagnates---reward hacking or gaming \citep{amodei2016concrete,skalse2022reward}---and over-optimisation of a learned reward model follows predictable scaling laws \citep{gao2023scaling}. Simulators of recommender systems inherit the biases of the logs they are fitted to \citep{huang2020keeping}, and offline evaluation is a recognised obstacle to deploying RL in the real world \citep{dulacarnold2021challenges}. KT--RL recommenders sit at the intersection of these problems because one learned model serves as dynamics, reward and evaluation metric. From epidemiology we borrow the negative control, an analysis in which the true effect is known to be absent so that any estimated effect reveals bias \citep{lipsitch2010negative,arnold2016negative}.

Our work connects these threads for programming exercise recommendation. We fit six learner models of equal predictive quality to the same real data and show how their disagreement propagates to policy rankings, and we add two instruments that the robust-evaluation literature has not used: a no-learning negative-control world, in which the true gain of every policy is exactly zero, and an exact decomposition of a tracker's reported gain into drift and surprise, with formal results that identify in-simulator gain with accumulated drift.
